\documentclass[12pt,a4paper]{article}

% -------- PACKAGES (Engineering Standard) --------
\usepackage[utf8]{inputenc}
\usepackage[T1]{fontenc}
\usepackage[a4paper, margin=1in]{geometry}
\usepackage{titlesec}
\usepackage{xcolor}
\usepackage{hyperref}
\usepackage{longtable}
\usepackage{booktabs}
\usepackage{enumitem}
\usepackage{array}
\usepackage{float}
\usepackage{graphicx}
\usepackage{listings}
\usepackage{fancyhdr}
\usepackage{tikz}

% -------- BRANDING & COLORS --------
\definecolor{mainblue}{RGB}{0,102,204}
\definecolor{darkgray}{RGB}{50,50,50}
\definecolor{codebg}{RGB}{248,248,248}

% -------- SECTION STYLING --------
\titleformat{\section}{\Large\bfseries\color{mainblue}}{\thesection}{1em}{}[\titlerule]
\titleformat{\subsection}{\large\bfseries\color{darkgray}}{\thesubsection}{1em}{}

% -------- HEADER/FOOTER --------
\pagestyle{fancy}
\fancyhf{}
\lhead{\footnotesize SmartShop: Technical Submission}
\rhead{\footnotesize LUWI Mafuleti}
\cfoot{\thepage}

% -------- CODE BLOCK STYLE --------
\lstset{
    backgroundcolor=\color{codebg},
    basicstyle=\ttfamily\small,
    breaklines=true,
    frame=single,
    keywordstyle=\color{blue},
    commentstyle=\color{gray},
    language=PHP,
    showstringspaces=false
}

\hypersetup{colorlinks=true, linkcolor=mainblue, urlcolor=blue}

\begin{document}

% ============================
% COVER PAGE
% ============================
\begin{titlepage}
    \begin{tikzpicture}[remember picture, overlay]
        % Full-page background image
        \node at (current page.center) 
            {\includegraphics[width=\paperwidth,height=\paperheight]{LLDP-and-CDP-2.jpg}};
        
        % Semi-transparent white box for readability
        \fill[white, opacity=0.85] 
            (current page.north west) rectangle (current page.south east);
    \end{tikzpicture}

    \begin{center}
        % Logo at top
        \vspace*{1.5cm}
        \includegraphics[width=0.85\textwidth,height=4cm]{FST Marrakech.png} \\[1cm]

        % Main Technical Title
        {\Huge \textbf{FINAL MODULE PROJECT REPORT}} \\[0.5cm]
        {\LARGE \textit{``Back-End Web Development''}} \\[1cm]
        
        {\Large \textbf{Architectural Design, Hardening, and Production Deployment of the}} \\[0.3cm]
        {\Huge \textcolor{mainblue}{\textbf{SmartShop Platform}}} \\[1.5cm]

        % Authors
        {\large \textbf{Developed by:}} \\[0.3cm]
        {\Large MAFULETI Luwi} \\[2cm]

        % Supervisor
        {\large \textbf{Supervised by:}} \\[0.3cm]
        {\Large Prof. SADQI Yassine} \\[2.5cm]

        % Academic year
        {\large Academic Year 2025-2026}

        \vfill
        \textbf{Project Repository:} \href{https://github.com/yourusername/e-commerce-platform}{GitHub: SmartShop Platform}
    \end{center}
\end{titlepage}
\vfill
\begin{center}
    \textbf{Project Repository:} \href{https://github.com/yourusername/e-commerce-platform}{GitHub: SmartShop Platform}
\end{center}
\newpage

\tableofcontents
\newpage

% ==================================================
% PHASE I: CONCEPTUAL GENESIS & PLANNING
% ==================================================
\section{Phase I: Conceptual Genesis \& Architectural Planning}

In the initial stage of this project, my focus was on defining a robust architectural foundation. I aimed to transition from a traditional CRUD mindset to a scalable transactional system. I began by conducting a comprehensive feasibility study and architectural mapping, which I have archived in \texttt{LUWI\_Mafuleti\_IRISI\_Laravel\_Project\_Plan.pdf} and \texttt{LUWI\_Mafuleti\_IRISI\_Laravel\_Project\_Execution-3.pdf}. These documents served as my primary research guides for the initial MCD and MLD design.

\subsection{1. My Vision: Engineering for Scalability}
The objective was to engineer a high-integrity e-commerce engine. I defined three core pillars for the platform:
\begin{itemize}
    \item \textbf{High-Integrity Transactions:} Ensuring data consistency during the order lifecycle.
    \item \textbf{Actor-Based Access:} I implemented a granular separation between Guests, Members, and Administrators to protect the system's perimeter.
    \item \textbf{Relational Excellence:} I designed a normalized database structure that is fully extensible for future multi-partner integrations.
\end{itemize}

\subsection{2. System Requirements \& Stakeholder Mapping}
I mapped the ecosystem to ensure that each system actor had a defined boundary of interaction. I planned a frictionless consumer journey from product discovery to secure checkout, while simultaneously designing a control layer for catalog governance.

\subsection{3. Conceptual Data Model (MCD)}
The blueprint of SmartShop began with the \textbf{MCD}. I established strict relational cardinalities to prevent data drift. My design accounted for core entities like Users and Products, alongside production-grade entities like Variants and Payments.

\begin{figure}[H]
    \centering
    % \includegraphics[width=\textwidth]{mcd.png}
    \caption{My Conceptual Data Model - Defined as the blueprint for system integrity.}
\end{figure}

\subsection{4. Complete Logical Data Model (MLD)}
I translated my conceptual vision into a 3rd Normal Form (3NF) Logical Model. The following entities represent the full implemented state of my database:

\subsubsection{4.1 Core Transactional Entities}
\begin{itemize}[noitemsep]
    \item \textbf{USERS:} The identity core (\underline{id}, name, email, password, role, timestamps).
    \item \textbf{CATEGORIES \& PRODUCTS:} The taxonomy core (\underline{id}, name, price, stock, \#category\_id, timestamps).
    \item \textbf{CART \& CART\_ITEMS:} Active persistence layer (\underline{id}, \#user\_id, \#product\_id, quantity).
    \item \textbf{ORDERS \& ORDER\_ITEMS:} Immutable records (\underline{id}, \#user\_id, total\_price, status, timestamps).
\end{itemize}

\subsubsection{4.2 Extended Scalability Layer}
To support production-grade functionality, I implemented:
\begin{itemize}[noitemsep]
    \item \textbf{PRODUCT\_IMAGES \& VARIANTS:} Handling multi-media and SKU-level data (\underline{id}, \#product\_id, url, sku, price).
    \item \textbf{REVIEWS \& PAYMENTS:} Enabling user feedback and financial persistence (\underline{id}, \#user\_id, \#order\_id, amount, rating).
    \item \textbf{ADDRESSES:} Supporting residential delivery with primary-flag logic (\underline{id}, \#user\_id, city, country, is\_primary).
    \item \textbf{VENDORS \& VENDOR\_PRODUCTS:} Prepared the schema for a multi-seller marketplace (\underline{id}, name, \#product\_id).
\end{itemize}

\subsection{5. Hybrid Data Strategy}
I adopted a hybrid approach for data initialization. I first imported a verified \texttt{database\_dump.sql} to preserve development-time data and then utilized Laravel's migration system to ensure schema synchronization. This strategy guaranteed that my production environment was both data-rich and architecturally aligned.

\vfill
\begin{center}
    \textit{Further details on the initial design constraints and technical business rules can be found in my Project Plan and Execution PDFs.}
\end{center}

\newpage
% ==================================================
% PHASE II: TECHNICAL IMPLEMENTATION & EXECUTION
% Source: EXECUTION_LOG.tex, MASTER_DOCUMENTATION.tex
% ==================================================
\section{Phase II: Technical Implementation \& Execution}

With the architectural blueprint finalized, I transitioned into the execution phase. I adopted an iterative development approach, documented meticulously in \texttt{EXECUTION\_LOG.tex}. My primary focus was to maintain a clean separation of concerns while ensuring that the backend logic remained performant and secure.

\subsection{1. The MVC Engine: Model-View-Controller Mastery}
I implemented the core logic using Laravel’s MVC pattern. My objective was to keep controllers lean by offloading data logic to the Eloquent models and request validation to dedicated FormRequest classes.

\subsubsection{1.1 Eloquent ORM \& Relationship Orchestration}
I leveraged Eloquent to handle complex relational data. To prevent the common "N+1" performance issue, I enforced eager loading (\texttt{with()}) across all catalog queries.

\textbf{Performance Audit: My N+1 Query Resolution} \\
During my technical audit, I identified that the \texttt{ViewController} was executing individual queries for every product's category and images. I refactored the retrieval logic as follows:
\begin{lstlisting}
// BEFORE: Individual queries triggered during loop (Performance Leak)
// $products = Product::paginate(12);

// AFTER: Eager loading optimized for production (Linear Efficiency)
$products = Product::with(['category', 'images'])->paginate(12);
\end{lstlisting}

\textbf{Example: My Image Resolution Logic} \\
In the \texttt{Product} model, I developed a custom attribute to handle image fallbacks and standardized storage paths, ensuring that both local and remote (Unsplash) assets are resolved seamlessly:
\begin{lstlisting}
public function getImageUrlAttribute() {
    $firstImage = $this->images->first();
    if (!$firstImage || empty($firstImage->url)) {
        return 'https://images.unsplash.com/fallback-image-url';
    }
    return str_starts_with($firstImage->url, 'http') 
        ? $firstImage->url 
        : asset('storage/' . str_replace('storage/', '', $firstImage->url));
}
\end{lstlisting}

\subsection{2. Transactional Integrity: The Cart \& Checkout Lifecycle}
I engineered the checkout process to be atomic. I used database transactions to ensure that the user’s order is only finalized if the inventory is successfully decremented and the cart is cleared in a single, uninterrupted operation.

\begin{lstlisting}
DB::transaction(function () use ($cart) {
    $order = Order::create([
        'user_id' => auth()->id(),
        'total_price' => $cart->total_price,
        'status' => 'pending'
    ]);

    foreach ($cart->items as $item) {
        // I ensured stock is decremented immediately to prevent overselling
        $item->product->decrement('stock', $item->quantity);
        OrderItem::create([
            'order_id' => $order->id,
            'product_id' => $item->product_id,
            'quantity' => $item->quantity,
            'price' => $item->product->price
        ]);
    }
    $cart->items()->delete();
});
\end{lstlisting}

\subsection{3. Frontend Layer: Blade Componentization}
My UI strategy focused on modularity. I developed a suite of reusable Blade components (e.g., \texttt{product-card}) and a master layout (\texttt{app.blade.php}) to ensure design consistency. I also integrated a dynamic theme engine, allowing users to toggle between Light and Dark modes, with state persisted via \texttt{localStorage}.

\subsection{4. Background Task Orchestration}
To optimize the user experience, I offloaded time-consuming tasks like mailable generation to a database queue. This ensures that the user receives an immediate confirmation on the frontend while the system handles email delivery asynchronously via a dedicated background worker.

\vfill
\begin{center}
    \textit{The complete record of my implementation commands and versioned milestones is preserved in the MASTER\_DOCUMENTATION.tex file.}
\end{center}


\newpage
% ==================================================
% PHASE III: HARDENING, SECURITY & STABILIZATION
% Source: HARDENING_PLAN.tex, SESSION_HANDOVER.md
% ==================================================
\section{Phase III: Strategic Hardening \& System Stabilization}

The final phase of development was dedicated to transforming a functional prototype into a hardened, production-ready platform. I executed a series of strategic missions to secure administrative boundaries and ensure the long-term integrity of user data. My comprehensive strategy for this phase is outlined in \texttt{HARDENING\_PLAN.tex}.

\subsection{Mission: The Executive Command Center (Governance)}
I refactored the administrative experience into a centralized "Executive Command Center." My objective was to provide administrators with actionable intelligence rather than raw data.

\begin{itemize}
    \item \textbf{Real-Time Analytics:} I implemented a controller-driven dashboard that aggregates platform health metrics, including total revenue and active order counts.
    \item \textbf{Mission: Inventory Intelligence (Low Stock Pulse):} I developed an automated alerting system that highlights products with less than 5 units remaining, enabling proactive inventory management.
    \item \textbf{Member Moderation:} I built an administrative suite to manage the community, allowing for role elevation (User to Admin) and account governance.
\end{itemize}

\subsection{Mission: Defensive Perimeter (Security Hardening)}
I identified and mitigated several critical security risks during my stabilization audit:

\subsubsection{Mission: Identity \& Anti-Spoofing Architecture}
I identified a vulnerability where user profiles could potentially be accessed or modified via ID parameters in the URL. I re-engineered the routing architecture to be static, relying exclusively on \texttt{Auth::user()} in the controller.
\begin{lstlisting}
// I refactored the route to remove user IDs, ensuring absolute isolation
Route::put('/profile/update', [UserController::class, 'updateProfile'])
    ->name('profile.update');
\end{lstlisting}

\subsubsection{Mission: Concurrency Control (Atomic Inventory Security)}
A significant production risk I identified was the potential for "Overselling" during simultaneous checkouts. To ensure 100\% stock integrity, I implemented \texttt{lockForUpdate()} within my checkout transactions. This ensures that the system locks the product row during the stock check, preventing other requests from modifying it until the decrement is complete.
\begin{lstlisting}
// Ensuring stock integrity under high load
$product = Product::where('id', $item->product_id)->lockForUpdate()->first();
if ($product->stock >= $item->quantity) {
    $product->decrement('stock', $item->quantity);
}
\end{lstlisting}

\subsubsection{Mission: Data Integrity \& CSRF Protection}
I enforced strict \texttt{\$fillable} arrays across all models—most notably the \texttt{Address} and \texttt{User} models—to prevent malicious attribute injection. Furthermore, I ensured that every form submission is protected by Laravel's CSRF (\textbf{Cross-Site Request Forgery}) middleware. I also implemented dedicated \textbf{FormRequest} classes to decouple validation logic from my controllers, ensuring that only sanitized data enters the domain layer.

\subsection{Mission: Universal Accessibility (Mobile Excellence)}
I conducted a final audit of the user experience to ensure a premium feel across all devices.
\begin{itemize}
    \item \textbf{Mobile Navigation Overhaul:} I redesigned the navigation architecture to prevent UI overlapping on small viewports, ensuring 100\% accessibility for mobile shoppers.
    \item \textbf{Mission: The Messenger (Feedback):} I integrated a custom toast notification system to provide semantic feedback (Success/Error) during critical user actions.
\end{itemize}

\vfill
\begin{center}
    \textit{Detailed technical specifications on these missions are available in the HARDENING\_PLAN.tex and recent session handover notes.}
\end{center}


\newpage
% ==================================================
% PHASE IV: PRODUCTION DEPLOYMENT & INFRASTRUCTURE
% Source: DEPLOYMENT_GUIDE.tex
% ==================================================
\section{Phase IV: Production Deployment \& Infrastructure Orchestration}

The transition from local development to a live production environment was a multi-stage engineering challenge. I successfully deployed the platform on a hardened Linux environment, managing everything from infrastructure provisioning to server-level security negotiations. My full deployment protocol is documented in \texttt{DEPLOYMENT\_GUIDE.tex}.

\subsection{1. Infrastructure Acquisition (GitHub Education)}
I initiated the deployment by applying for the \textbf{GitHub Student Developer Pack}. Following a 72-hour identity verification process, I successfully claimed \$200 in DigitalOcean credits. For security and verification reasons, I linked my payment information, which included a \$5 initial account charge to activate the production-grade resources.

\subsection{2. Server Architecture: Choosing the Droplet}
During the architectural selection process, I evaluated the DigitalOcean App Platform but identified a critical constraint: it defaults to PostgreSQL. Given that SmartShop is optimized for \textbf{MySQL}, I chose to provision a custom \textbf{Droplet} (Ubuntu 24.04 LTS) to maintain full control over the environment.

\subsubsection{2.1 Mission: OOM Recovery \& Resource Stability}
During the initial MySQL installation on the 1GB RAM Droplet, I encountered several \textbf{Out of Memory (OOM)} crashes. To stabilize the environment and ensure high-availability for the database, I implemented a \textbf{2GB Swap file}. This allowed the system to offload inactive memory to the disk, preventing further crashes during peak resource usage:
\begin{lstlisting}[language=bash]
# I stabilized the server by implementing Swap memory
sudo fallocate -l 2G /swapfile
sudo chmod 600 /swapfile
sudo mkswap /swapfile
sudo swapon /swapfile
echo '/swapfile none swap sw 0 0' | sudo tee -a /etc/fstab
\end{lstlisting}

\subsubsection{2.2 Database Migration \& LEMP Stack}
I followed a ``Database-First'' deployment strategy. I first initialized a secured MySQL instance on the server, imported my \texttt{database\_dump.sql}, and then configured the Droplet to merge the application files via Git. This hybrid strategy ensured that my development-time data and schema were perfectly synchronized in production.

\subsection{3. Mission: The Mailing & Persistence Resolution}
A significant production hurdle was the default security policy of DigitalOcean, which blocks outgoing ports \textbf{465 and 587}. I resolved this by engaging with DigitalOcean Support, successfully unblocking the ports for real-time notifications. 

To ensure the persistence of my background workers, I implemented a \textbf{Supervisor} configuration to manage the Laravel queue. This guarantees that my mailable processes are automatically restarted if they crash or following a server reboot:
\begin{lstlisting}[language=bash]
# I ensured mailable persistence with Supervisor
[program:smartshop-worker]
command=php /var/www/smartshop/artisan queue:work --sleep=3 --tries=3
autostart=true
autorestart=true
user=www-data
\end{lstlisting}

\vfill
\begin{center}
    \textit{The complete record of my infrastructure configuration, Nginx setup, and SSL encryption (Certbot) is preserved in the DEPLOYMENT\_GUIDE.tex file.}
\end{center}

% ==================================================
\newpage
\section{V. My Visual Evidence (System Demonstration)}

I have compiled a comprehensive gallery of the live system (\texttt{smartshop-luwi.tech}) to demonstrate the platform’s functionality and my version control methodology.

\subsection{1. Core Storefront Journey}
\begin{figure}[H]
    \centering
    % \includegraphics[width=\textwidth]{homepage.png}
    \caption{Homepage - The entry point to the Discovery Engine.}
\end{figure}

\begin{figure}[H]
    \centering
    % \includegraphics[width=\textwidth]{collections.png}
    \caption{Collections (Shop) - Advanced filtering and search in action.}
\end{figure}

\subsection{2. User Identity \& Commerce Flows}
\begin{figure}[H]
    \centering
    % \includegraphics[width=\textwidth]{profile.png}
    \caption{Member Profile - Secure identity and residential management.}
\end{figure}

\begin{figure}[H]
    \centering
    % \includegraphics[width=\textwidth]{cart.png}
    \caption{Persistent Cart - Real-time total calculation and item management.}
    \centering
    % \includegraphics[width=\textwidth]{orders.png}
    \caption{Order History - Tracking the full transactional lifecycle.}
\end{figure}

\subsection{3. Administrative Command Center}
\begin{figure}[H]
    \centering
    % \includegraphics[width=\textwidth]{admin.png}
    \caption{Admin Dashboard - Centralized analytics and management.}
\end{figure}

\subsection{4. Git Methodology (Version Control Evidence)}
\begin{figure}[H]
    \centering
    % \includegraphics[width=\textwidth]{git_history.png}
    \caption{Git Commit History - Documenting my atomic, feature-branch development workflow.}
\end{figure}

\vfill
\begin{center}
    \textit{Additional screenshots for Discover, Story, Support, Login, and Signup are included in the final technical archive submitted via email.}
\end{center}

% ==================================================
\chapter{Conclusion: Reflections on a Production Journey}

The development and deployment of the \textbf{SmartShop} platform have been a definitive exercise in transitioning from conceptual software design to a hardened, production-grade reality. This project allowed me to demonstrate mastery over the Laravel ecosystem, not merely as a CRUD framework, but as a robust engine for high-integrity transactional commerce.

Throughout this journey, I prioritized three core engineering principles:
\begin{itemize}
    \item \textbf{Architectural Integrity:} By enforcing a strict 3NF database schema and a layered MVC structure, I built a foundation that is technically prepared for future multi-partner scaling.
    \item \textbf{Resilient Infrastructure:} Navigating the complexities of DigitalOcean—from OOM recovery via Swap memory to negotiating server-level SMTP port security—proved my ability to manage the full lifecycle of a production system.
    \item \textbf{Uncompromising Security:} From atomic concurrency control to anti-spoofing route refactoring, I ensured that system integrity remains the perimeter that protects both the user data and the transactional record.
\end{itemize}

SmartShop stands as a mature, stable, and highly performant platform. It represents the successful convergence of technical optimization and professional deployment, ready to handle real-world traffic while maintaining a premium user experience. This project has solidified my expertise in back-end engineering, infrastructure orchestration, and the strategic hardening required for modern web applications.

\end{document}

